% ACS Nano / Nano Letters Template
% Change journal= to: nalefd (Nano Letters), ancac3 (ACS Nano)
\documentclass[journal=nalefd,manuscript=article]{achemso}

\usepackage{amsmath,amssymb,bm}
\usepackage{graphicx}
\usepackage{xcolor}
\usepackage{hyperref}
\usepackage{siunitx}
\usepackage{booktabs}
\usepackage[version=4]{mhchem}

% Placeholder for missing figures
\newcommand{\placeholder}[1]{%
  \begin{center}
    \fbox{\parbox{0.9\columnwidth}{\centering\small\textit{[FIGURE: #1]}}}
  \end{center}
}

% ─── Authors ─────────────────────────────────────────────────────────────────
\author{Author One}
\affiliation{Department of Applied Physics, University, Country}

\author{Author Two}
\affiliation{Department of Applied Physics, University, Country}

\author{Corresponding Author}
\email{corresponding@email.com}
\affiliation{Department of Applied Physics, University, Country}

% ─── Title ───────────────────────────────────────────────────────────────────
\title{Title of the Paper}

% ─── Abstract ────────────────────────────────────────────────────────────────
\begin{document}

\begin{abstract}
Abstract text here (≤150 words). First sentence states the key result.
Context sentence. Problem sentence. Approach sentence. Key result with numbers.
Significance/implication sentence.
\end{abstract}

\begin{tocentry}
  % Table of contents image (ACS requires this)
  \placeholder{TOC graphic: visual summary of the work}
\end{tocentry}

% ─── Introduction ────────────────────────────────────────────────────────────
\section{Introduction}

Broad context and motivation.\cite{ref1}
Specific problem and gap in knowledge.\cite{ref2,ref3}
Prior approaches and limitations.\cite{ref4}
Here we demonstrate...\cite{ref5}

% ─── Results and Discussion ──────────────────────────────────────────────────
\section{Results and Discussion}

\subsection{Device Architecture and Characterization}

\begin{figure}[htbp]
  \centering
  \placeholder{Device schematic and characterization data (AFM, Raman, MOKE)}
  \caption{\textbf{Device architecture.}
           (a) Schematic cross-section.
           (b) AFM topography.
           (c) Raman spectra.
           (d) MOKE hysteresis loop at \SI{300}{\kelvin}.}
  \label{fig:device}
\end{figure}

\subsection{Main Result Subsection Title}

\begin{figure*}[htbp]
  \centering
  \placeholder{Main result figure spanning two columns}
  \caption{\textbf{Main results.}
           (a)--(d) Description of panels.}
  \label{fig:main}
\end{figure*}

% ─── Conclusion ──────────────────────────────────────────────────────────────
\section{Conclusion}

In summary, we have demonstrated...
The key result of X shows...
Our results open a route toward...

% ─── Methods ─────────────────────────────────────────────────────────────────
\section{Experimental Methods}

\subsection{Sample Preparation}

\subsection{Optical Measurements}

\subsection{Data Analysis}

% ─── Acknowledgments ─────────────────────────────────────────────────────────
\begin{acknowledgement}
The authors acknowledge...
This work was supported by...
\end{acknowledgement}

% ─── Supporting Information ──────────────────────────────────────────────────
\begin{suppinfo}
Figure S1: [Description].
This material is available free of charge via the Internet at \url{http://pubs.acs.org}.
\end{suppinfo}

% ─── References ──────────────────────────────────────────────────────────────
\bibliography{refs}

\end{document}
